\section{A quartic presentation with 394 variables}
\label{sec:v13}

The transformations in this section give the following specified
polynomials. The arithmetic inputs and their scope are those retained
in Section~\ref{subsec:v12-inputs}.
\begin{theorem}\label{thm:v13-quartic}
There is a specified integer polynomial $P_{394}$ of total degree four
in 394 variables such that
\[
 X\in\mathbb Z\quad\Longleftrightarrow\quad
 \exists\mathbf v\in\mathbb Q^{393}:P_{394}(X,\mathbf v)=0.
\]
It has 6,338 nonzero monomials, content one, and largest absolute
coefficient $16\cdot3^{2^{26}}$. Its symbolic coefficient table uses
$C=3^{2^{25}}$ as a fixed integer, not an additional variable.
There is also a specified integer quartic $P_{411}$ with
\[
 X\in\mathbb Z\quad\Longleftrightarrow\quad
 \exists\mathbf w\in\mathbb Q^{410}:P_{411}(X,\mathbf w)=0.
\]
It has 6,388 nonzero monomials, content one, and all coefficients have
absolute value less than $2^{3247}$. Its largest coefficient has exactly
3,247 bits; its exact value is included in the coefficient artifact.
\end{theorem}

Some transformations below are inverse rational changes of private
witnesses, fixing every outside coordinate. Others change the arithmetic
certificate or enlarge a point domain; for those we prove soundness and
existence of witnesses for every integer input separately.

\subsection{Conventions and retained arithmetic inputs}
Put $K=\mathbb Q(\rho)$, $\rho^2=2$, $d_0=3-\rho$, $D=d_0^2$, and
$E:y^2=x^3-Dx$. Retain $2E(K)=\mathbb ZP$, with
\[
 P=(c_*,a_*b_*),\quad a_*=-1+7\rho/4,\quad
 b_*=39/8-9\rho/2,\quad c_*=a_*^2.
\]
The roots $Q_i=[n_i]P$ have nonzero odd indices, with
$(n_1,n_2,n_3,n_4)=(n,a,t,b)$. In the intermediate source the auxiliary point is
$H=[h]P\ne O$, and $U_i=[u_i]P$, $V_0=H-P$, $V_i=U_i-Q_i$.
We will allow $U_i$ and $V_i$ to be $O$ for $i>0$. The proof below
still excludes $V_0=O$. The scalar point is $R=Q_2-P=[s]P$, where
$s=a-1$ is nonzero and even.

The retained root-height implication is
\begin{equation}\label{eq:v13-height}
 B^2-6<7\mathcal H|N|^{19/k},\qquad
 B=\max_i|n_i|,\quad 1\le\mathcal H<2^{2093966},
\end{equation}
whenever the $k$th powers of the three contact radicals divide the
nonzero integer $N$. Their scales are $1,C_0^2,(1+C_1^2)^2$, where
$\prod_i x(Q_i)=C_0+C_1\rho$. The retained scalar anchor gives
\begin{equation}\label{eq:v13-anchor}
 \log_2|s|>25165823.5.
\end{equation}
These are precisely the root-height and anchor statements proved earlier;
neither is inferred from a finite coefficient calculation.

\subsection{Rational logarithms and the two order-ten tests}
Write $f=1/x$, $z=-x/y$, and extend these regular functions by
$f(O)=z(O)=0$. Define
\begin{align*}
 Q(w)&=1-\frac{1526}{1073}w+\frac{513}{1073}w^2,\\
 N(w)&=1-\frac{9776}{5365}w+\frac{2941}{3219}w^2
 -\frac{1024}{13949}w^3-\frac{1920}{237133}w^4
 -\frac{1024}{711399}w^5-\frac{1408}{5928325}w^6
 +\frac{2432}{103152855}w^8,\\
 \mathcal R(w)&=N(w)/Q(w).
\end{align*}
The exact identity
\[
 QG_9-N=\frac{2970726}{254443709}w^{10}
       -\frac{41697666}{6361092725}w^{11}
\]
proves that this replacement has the same order-forty local accuracy.
Indeed, on $[0,1]$ the denominator is at least $60/1073$ and the
Bernstein coefficients of the numerator are
\[
 1,\ \frac{4143}{5365},\ \frac{260069}{450660},\
 \frac{27829}{67340},\ \frac{4636683}{16599310},\
 \frac{82391}{474266},\ \frac{2245591}{23713300},\
 \frac{955437}{23713300},\ \frac{83904}{10113025}.
\]
They are all positive and give $N\ge83904/10113025$. Every nonidentity point of $2E(K)$ has
$0<Df^2<1$ at both real embeddings; at $O$ its value is zero.
Thus it is defined and nonzero on these doubled-point domains. For
the later enlarged H domain and its difference H-P, we instead use
the global numerator and denominator nonvanishing certificates below.

At an unramified place above $p>76176$, a nonidentity formal point of
depth $r=v(z)\ge1$ has $v(Df^2)=4r$. All coefficient denominators
above are supported at primes at most37. Consequently
$v(\mathcal R-G_9)\ge40r$. The earlier full-tail estimate gives
$v(\log_E(Q)-z\mathcal R(Df^2))\ge41r$.
The point $O$ has exact zero error. We retain the $G_2$ numerator at
$V_1$, of relative precision12, and use $\mathcal R$ at $V_0,V_2,V_3,V_4$.

Set $\ell_0=z(V_0)\mathcal R(Df(V_0)^2)$ and define
\[
 L_i=z(U_i)/z(H),\quad
 L'_1=\overline{z(V_1)G_2(Df(V_1)^2)/\ell_0},\quad
 L'_i=\overline{z(V_i)\mathcal R(Df(V_i)^2)/\ell_0}\quad(i>1).
\]
The nonzero denominator is proved below, before these quotients are used.
The tests have orders1 on $f(U_i)$ and $\overline{f(V_i)}$, order40
on $f(H)$, order10 on $L_1-L'_1,L_1-\beta$, and order40 on the other
three cross-tests and
\[
 F(L_2,L_3,L_4),\qquad F(A,T,B)=T^2-A(A-2)TB+(A-2)^2B^2-1.
\]
Conjugation and the place conventions are unchanged.

For completeness, spell out the zero-safe local argument. Let $A(Q)$ be
the analytic unit defined by $\log_E(Q)=z(Q)A(Q)$, with $A(O)=1$.
Define $T_2=G_2/A$ and $T_{40}=\mathcal R/A$ by their power series.
These units equal1 at $O$. Their errors at nonidentity formal points
have the stated orders. The primary ratio is exactly
$(u_i/h)A(H)/A(U_i)$, also when $u_i=0$. The alternate ratio is
$(u_i-n_i)/(h-1)$ times the corresponding conjugated unit ratio.
At least one of $h,h-1$ is a $p$-unit. Therefore one full ratio is
integral, and the positive-order cross-test transfers integrality to
both approximate and full ratios. Only now infer depth at least20 for
a nonzero $U_i$; if $U_i=O$, its error is zero. The primary error has
order80. Multiplication by $h,h-1$ and subtraction give
\[
 L_1\equiv n\pmod{p^{10}},\qquad L_i\equiv n_i\pmod{p^{40}} (i>1).
\]
No quotient $\log_E(O)/z(O)$ is used.

\subsection{Odd roots, shared reciprocals and identity-inclusive charts}
The retained finite chart for a root is
$x=h_Q^2$, $j_Q^2=x^2-D$, $Q=(x,h_Qj_Q)$.
It has $h_Qj_Q\ne0$ by the real embeddings and $\operatorname{Norm}(d_0)=7$.
Replace its old odd-index witnesses by $A\in K,u\in\mathbb Q$ with
\[
 A^2-2(x+j_Q)A-D=0,\qquad \operatorname{Norm}(A)=-u^2.
\]
Here is the reverse implication, including the squareclass choice.
Put $q=x+j_Q$ and $t=(A-q)/h_Q$. The equation gives
$t^2=2q$. The point $(A,At)$ lies on $E$ and its tangent slope is
$h_Q+t$, so its double is exactly $(x,h_Qj_Q)$.
The inherited squareclass decomposition writes its abscissa class as
$[\epsilon]^m[-1]^e[d_0]^f$, where $\epsilon=1+\rho$ and $e,f\in\{0,1\}$.
The norm condition forces $m$ odd and $f=0$. Thus $[A]=[\epsilon]$
or $[-\epsilon]$. In the latter case replace $A,t$ by $2q-A,-t$;
the product of the two abscissae is $-D$, so the new class is
$[\epsilon]$. Choosing the corresponding square root recovers the old
odd witnesses without changing the root. Conversely the old witnesses
$t,r$ give $A=th_Q+q=\epsilon r^2$, $u=\operatorname{Norm}(r)$.
This is equality of exterior projections, not a unique private inverse.

For the primary ratios keep one shared field value
$S=j_H/h_H=-1/z(H)$. The finite-point equations become
\[
 h_H S=j_H,\qquad j_iL_i=h_iS.
\]
These remove the four private primary parameters $z(U_i)$ without
changing their ratios. The scalar equation can be fused in the same way.
If its old equations are $Tz_R+a_*N=0$ and
$h_H\varrho+j_Hz_R=0$, replace them by
\[
 T\varrho-a_*NS=0.
\]
The inverse is $z_R=-\varrho/S$; $S\ne0$ follows from the retained
finite H chart.

For $i>0$ an identity-inclusive primary chart instead uses
\[
 f_i=r_i^2,\qquad J_i^2+Df_i^2=1,\qquad J_iL_i=r_iS.
\]
When $r_i\ne0$, set $h_i=1/r_i,j_i=J_i/f_i$ to recover the old
finite chart. When $r_i=0$, interpret $U_i=O$. Both choices $J_i=\pm1$
give that same point. The factor $J_i$ cannot vanish: it would give
$D r_i^4=1$, and taking real signs and then norms gives
$7\operatorname{Norm}(r_i)^2=1$, impossible over $\mathbb Q$.
Thus $z(U_i)=-r_i/J_i$ and the primary ratio is correct also at $O$.

Write the corresponding fixed root as $Q_i=(c,ab)$ with
$c=a^2$, $b^2=c^2-D$. Define
\[
 d=c+Df_i,\quad dg=aJ_i-br_i,\quad u=cr_i+bg,
 \quad dk=bJ_i-2au,\quad f(V_i)=g^2,\quad kz(V_i)=-g.
\]
The factor $d$ is positive at both real embeddings. Reduction by the
primary conic gives $k^2+Dg^4=1$, so $k\ne0$ by the same norm argument.
On the ordinary subtraction domain these expressions give
$V_i=(1/g^2,k/g^3)=U_i-Q_i$. They have the following exact boundary values.
If $r_i=0$, then $g=J_i/a$, $k=-bJ_i/c$, and $V_i=-Q_i$.
If $g=0$, substitution gives
$(cr_i^2-1)(c+Dr_i^2)=0$, hence $U_i=Q_i$ and $V_i=O$.
The remaining equal-abscissa case is the regular doubling limit for
$U_i=-Q_i$. Equality with the group-law morphism follows from its equality
on the dense ordinary domain and regularity of these formulas.

For H and the denominator difference retain the finite chart, or its
following exact rescaling. With $x=h_H^2$, put
\[
 d=c_*x+D,\quad N=a_*j_H-b_*h_H,\quad dg=N,
 \quad T=-2Dh_H-a_*b_*j_H+2Db_*g.
\]
The scalar reciprocal elimination is the identity
$L=a_*dZ-zK-a_*G$, where
$G=dg-N$, $K=a_*dk-T$, $Z=kz+g$, $L=Tz+a_*N$.
Since $a_*d\ne0$, reinserting $k=T/(a_*d)$ recovers the old equations.

For V0 one may also eliminate $j_H$ by writing $j_H=h_HS$ and $g=h_Hv$.
The resulting exact equations are
\[
 S^2=x-Df(H),\quad (c_*x+D)v=a_*S-b_*,\quad
 f(V_0)f(H)=v^2,
\]
\[
 (-2D-a_*b_*S+2Db_*v)\ell_0
       +a_*(a_*S-b_*)\mathcal R(Df(V_0)^2)=0.
\]
Both maps divide only by the already nonzero $h_H$ and $x$.
These exact maps precede the separately justified enlargement of H below.

\paragraph{The denominator cannot be the identity.}
The old argument using $V_1\ne O$ is no longer available. Instead suppose
$V_0=O$. The reciprocal chart gives $H=P$, and
\[
 f(H)=\frac{8(57+28\rho)}{1681}
\]
is a unit at every consumed place over $p>76176$. The one-place detector
lemma would force its order-forty test at any contact prime, a contradiction.
Thus all three contact radicals equal1. Equation~\eqref{eq:v13-height}
with $N=1$ gives $B^2-6<7\mathcal H$, contrary to
$B\ge|a|>(40\mathcal H)^{10}$ from the unchanged anchor.
This argument precedes every interpretation of an LP quotient.
Now $V_0\ne O$, its parameter is nonzero, and the established positivity
of $\mathcal R$ gives $\ell_0\ne0$. The local ratio proof above therefore
applies to all new solutions. It proves soundness directly, without
asserting that identity-branch witnesses lift to the old affine source.

\subsection{Scaled biquadratic comparison charts}
For one root write $Q=(c,ab)$, $c=a^2$ and $b^2=c^2-D$.
The reciprocal comparison chart above uses $r,J,f,g,k,u,F$, with
$f=r^2$, $F=g^2$. Its retained squared outputs are
$Y_U=f^2$ and $Y_V=F^2$, as elements of $K$, not separate component
squares. Replace the seven field values by five:
\[
 V=ar,\quad Q'=ag,\quad I=aJ,\quad K'=ak,\quad W=rg.
\]
The new equations are
\begin{align*}
 VQ'&=cW,& V^2+(Q')^2+DW^2+2bW&=1,\\
 aI&=bV+cQ'+DVW,& aK'&=-bQ'-cV-DQ'W,\\
 IL&=VS,& K'z&=-Q',\\
 DcY_U+I^2&=c,& DcY_V+(K')^2&=c.
\end{align*}
Every displayed equation is quadratic in retained coordinates. In
particular c is an existing coordinate, so cY is a quadratic product.

Here is an algebraic proof including all zero branches. Put
$d=c+Dr^2$, $e=c+Dg^2$, and define
\[
 T=c(r^2+g^2)+Dr^2g^2+2brg-1,
\]
\[
 G=dg-aJ+br,\quad E'=ak+bg+er,\quad
 H'=dk-bJ+2a(cr+bg).
\]
Modulo only $c=a^2$ and $b^2=c^2-D$, direct expansion gives
\begin{align*}
 a^2(J^2+Dr^4-1)&=dT-G(dg+br+aJ),\\
 a^2(k^2+Dg^4-1)&=eT+E'(ak-bg-er),\\
 dE'&=aH'+DrT-bG.
\end{align*}
The factors a,c,D are nonzero, and d,e are positive at both real
embeddings. An old solution has G=H'=0, hence T=E'=0, proving
all new equations and the second reciprocal conic. Conversely set
$r=V/a$, $g=Q'/a$, $J=I/a$, $k=K'/a$; the product equation
forces W=rg. The new equations give T=G=E'=0, so the three
identities recover both conics and $H'=0$. Set $f=r^2$, $F=g^2$ and
$u=cr+bg$. Comparing each new Y equation with its recovered conic
gives $Y_U=r^4$ and $Y_V=g^4$. Dividing the ratio equations by a
restores the old ratios. Thus these are inverse rational maps,
fixing every outside coordinate.

There is no division by r,g,J,k,W,L,z or either test value. When
r=0 the first point is O and its ratio is zero; when g=0 the
difference is O and its parameter is zero. The biquadratic equation
prevents both from vanishing. These are exactly the already accepted
identity branches. This transformation replaces the intermediate inverse-addition change;
it is not composed with a second copy of that elimination.

\subsection{Enlarging the auxiliary H point domain}
After the exact H.j rescaling, the two coordinates H.h occur only
in the two component equations x=H.h$^2$. Delete them and keep
\[
                  xf=1,\qquad S^2=x-Df.
\]
They define the finite point H=(x,xS) of E(K), with x nonzero,
but no longer require H to be a double. The rank and torsion
calculation in the generator proof gives
\[
 H=[m]R_0+T,\quad P=2R_0,\quad m\in\mathbb Z,\quad T\in E[2](K).
\]
Write h=m/2 only for its rational free index. It is not notation
for multiplication by a half-integer on the elliptic curve.

We first exclude all finite points in $E[2]$, $P+E[2]$ and
$-P+E[2]$ before interpreting ratios. The points O and (0,0) are
excluded by the finite chart. Put $d_0=3-\rho$, $c_*=x(P)$.
The exact norm and translation identities are
\[
 \operatorname{Norm}(c_*)=41^2/64,\quad
 \operatorname{Norm}(c_*-d_0)=17^2/64,\quad
 \operatorname{Norm}(c_*+d_0)=63^2/64,
\]
\[
 x(P+(0,0))=-D/c_*,\quad
 x(P+(d_0,0))=d_0(c_*+d_0)/(c_*-d_0),
\]
\[
 x(P+(-d_0,0))=-d_0(c_*-d_0)/(c_*+d_0).
\]
Each of these abscissae and its inverse is integral outside
$\{2,3,7,17,41\}$. The same is true for the nonzero torsion
abscissae, and negating P changes none of these abscissae. Thus
f(H) is a unit at every consumed place. The retained order40 f(H)
test would then rule out every contact prime. All three radicals
would equal1, contradicting the root-height bound and scalar anchor
exactly as in the preceding V0 exclusion. This proof needs no
ratio involving H. Consequently m is neither0 nor2, H and H-P are
nontorsion, S is nonzero, and H is neither P nor -P.

In the rescaled difference equations put
\[
 d=c_*x+D,\quad N=a_*S-b_*,\quad dv=N,\quad
 T'=-2D-a_*b_*S+2Db_*v.
\]
The former positivity proof for d no longer applies. If d=0, the
curve equation gives $x=-D/c_*$, $S=\pm b_*/a_*$, so H is one of
the excluded points $\pm P+(0,0)$. If N=0, it gives
$(x-c_*)(x+D/c_*)=0$, also excluded. Thus d,N,v are nonzero.
Direct subtraction now gives
\[
 x(H-P)=1/(xv^2),\qquad y(H-P)=T'/(a_*dxv^3),
 \qquad z(H-P)=-a_*dv/T'.
\]
There is no unresolved equal-abscissa case, since $x=c_*$ means
H=P or -P. Since H-P is nontorsion, its ordinate and T' are
nonzero. In particular the retained equation $f(V_0)f(H)=v^2$ has
its intended interpretation, even if either real abscissa is negative.

The old positivity interval for the rational logarithm is no longer
enough. A stronger certificate is available: N(w) above has no root
modulo71, and Q(w) has no root modulo7. At the respective prime all
coefficients are integral and the leading coefficient is a unit.
These assertions are finite substitutions through all residues. Both
primes split in K, with square roots of2 equal to12,59 modulo71
and3,4 modulo7. A K-root would give a local root, necessarily integral,
and then a forbidden residue root. Hence N and Q never vanish in K.
The actual equation $T'\ell_0+a_*N\mathcal R(Df(V_0)^2)=0$
therefore implies $\ell_0\ne0$ and gives the correct denominator
logarithm. No sign assertion for this logarithm is needed.

The formal ratio argument changes by using integral group identities:
\[
 [m]U_i=[2u_i]H,\qquad
 [m-2]V_i=[2(u_i-n_i)]V_0.
\]
They kill T. Applying the formal logarithm only to formal points
gives ratios $u_i/h$ and $(u_i-n_i)/(h-1)$. At every consumed odd
prime h and h-1 are integral and differ by1. The earlier transfer
and subtraction proof therefore applies, including zero numerators.
At opposite split places use the same conjugation as before; at an
inert place it preserves the valuation.

The scalar transfer uses the same integral-index convention. For $R=[s]P$ use
\[
                    [m](2R)=[4s]H.
\]
If the scalar coefficient test has depth $d_s$, formal logarithm
isometry gives $v_p(s/(m/2))\ge d_s$. Since m/2 is p-integral,
$v_p(s)\ge d_s$. This includes the selected depth24. The proof
that 2R is formal when R is near O or (0,0) is unchanged. Thus
both the comparison and scalar height arguments survive the larger
domain. Completeness uses the original H in $2E(K)$ and forgets its
two square witnesses. It need not realize any newly admitted half-index
or torsion translate.

\subsection{A positive scalar input and the depth-twenty-four certificate}
Use the selected label and contact polynomial
\[
 \tau=2X+1,\qquad \beta=51\tau^2+2504,
 \qquad r=\beta(\beta-1)(\beta-2401)(32\beta+343)(2\beta-1).
\]
For every rational $X$, $\beta\ge2504$, hence $r>0$.
For integral $X$, $\tau$ is odd and $\beta$ is a positive odd integer.
Conversely, the final conclusion must be that $\beta$ is an \emph{odd}
integer. If $\tau=u/v$ in lowest terms, integrality gives $v^2\mid51$;
as $51$ is squarefree, $v=1$. Oddness of $\beta$ then makes $\tau$ odd
and $X=(\tau-1)/2$ integral. Mere integrality would not suffice:
$X=-1/2$ gives the even integer $\beta=2504$.
The unchanged odd chart for $Q_1=[n]P$ supplies the required oddness
once $n=\beta$ has been proved.

At a consumed prime $p>76176$, $51$ is a unit. A negative valuation
of $\beta$ is $2v_p(\tau)$, hence even. The five contact factors then
have the same valuation, so an odd valuation of $r$ must be positive.
This argument makes no assertion at the exceptional primes3 and17.

For integral $X$ the following simultaneous local calculation holds.
Since $\tau$ is odd, $\beta\equiv3\pmod8$; the two even contact
factors each have valuation one and equal odd parts modulo8. The
other factors have residues3,7,5. Therefore $v_2(r)=2$ and
$r/4\equiv1\pmod8$. At3, $\beta\equiv2$ and the first four factors
have product1 modulo3. Furthermore
\[
 (2\beta-1)/3=34\tau^2+1669\equiv1+\tau^2\pmod3,
\]
which is nonzero. Thus $v_3(r)=1$. At17, the contact residues are
5,4,1,10,9 with product15, a nonzero square. Thus $v_{17}(r)=0$.

An intermediate one-choice construction used $c=33033$ when
$3\mid\tau$ and $c=33$ otherwise. Both are square units at2 and17,
and their unit parts at3 are respectively1 and2, correcting the two
possible classes of $r$. The selected construction removes even
this choice and sets
\[
                 c=r+4896,\qquad4896=2^5\cdot3^2\cdot17.
\]
At2 the quotient $c/r=1+4896/r$ belongs to $1+8\mathbb Z_2$;
at3 and17 it belongs to $1+p\mathbb Z_p$. Hence $cr$ is a square
in all three local fields. Also $c,r$ are positive integers.
At a consumed scalar contact prime, positivity of the odd valuation
$v_p(r)$ was proved before using $c$; since4896 is a unit there,
$c$ is a unit there too. We do not assert that $c$ is a unit at
\emph{every} large prime. Its additional finite prime support is
included in the preparation before the continuity modulus is chosen.

The physical input square $Y=X^2$ is retained. In each other consumer,
the old occurrence of $Y$ is replaced by $102Y+102X+76$, since
$2(102Y+102X+76)+2403=51(2X+1)^2+2504$. This substitution must
include every scalar test and height-input consumer, not the defining
row $Y=X^2$ itself.

For the deterministic correction the two old physical rows, in their
fixed normalization, are
\[
 cs_0-20619812F=0,\qquad
 (7895415c+20619812F)V-20619812b_0=0.
\]
The retained contact expression has $r=-4(s_0+7895415)$.
Substitution and multiplication by $-4$ give
$rcV+82479248b_0=0$. With the already retained coordinate $r_2=r^2$,
the selected row is
\[
                    (r_2+4896r)V+82479248b_0=0.
\]
This is quadratic and preserves the physical sign of the scalar input.
The free correction coordinate and its product coordinate have no
remaining consumers. In the reverse arithmetic interpretation the
correction is defined to be $r+4896$; the removed finite-choice equation
is not reimposed.

Write $q=\operatorname{Norm}_{K/\mathbb Q}(f(H))$ and
$z(R)/z(H)=z_0+z_1\rho$. The selected scalar root is
\begin{equation}\label{eq:v13-scalar24}
 B_s=\frac{z_0^2+r z_1^2}{r^{32}},\qquad
 W=\frac{B_s^2+r q^2}{r^{32}},\qquad \mathfrak b=crW^2.
\end{equation}
Here $q\ne0$ and $r>0$, so $W>0$. Forward powers compute $r^{32}$;
there is no need to retain an inverse-power tower or $r^{64}$.
The square $q^2$ is computed without a separate norm coordinate:
if $f(H)=f_0+f_1\rho$, use
\[
 q^2=(f_0^2+2f_1^2)^2-8(f_0f_1)^2.
\]
The factor $8$ is essential: the second stored quantity is a product,
not the full $\rho$-coefficient.

Here is the changed local application. At a consumed prime let
$\eta=v_p(r)>0$ be odd. The correction is a unit there.
The pole obstruction gives $\eta+2v_p(W)>0$.
The two terms in each numerator of~\eqref{eq:v13-scalar24} have
opposite valuation parity, so they cannot cancel; assign valuation
$+\infty$ to a zero term. Successively this gives
\[
 v_p(B_s)>63\eta/4,\quad v_p(q)>61\eta/4,\quad
 v_p(z_0)>191\eta/8,\quad v_p(z_1)>187\eta/8.
\]
Thus both rational coefficient valuations are at least24.
Positive valuation of the norm supplies a place with $v(f(H))>0$;
there $H$ is formal and $v(z(R)/z(H))\ge24$.
Small $z(R)$ alone does not imply that $R$ is formal. If $e=v(x(R))$,
then $e\ne0$, and the duplication identity gives, for either sign of $e$,
\[
 v(x(2R))=-2v(z(R)),\qquad v(z(2R))=v(z(R)).
\]
For the enlarged H domain apply the formal logarithm to
$[m](2R)=[4s]H$, where $h=m/2$, to obtain
$v_p(s/(m/2))\ge24$. Since $m/2$ is p-integral at this odd prime,
$v_p(s)\ge24$. This uses integer group multiples only.
The same height theorem therefore gives
\[
 \operatorname h(\beta)\le\mathcal H|s|^{19/24}.
\]
The exact inequality
$(5/24)\,25165823.5>2093966+7$ and~\eqref{eq:v13-anchor} imply
$\operatorname h(\beta)<|s|/80\le|a|/40$.
The retained anchor also gives $\log(48\mathcal H)<|a|/40$.

For completeness choose $c$ first. The existing preparation permits
arbitrary finite coefficient-depth conditions over the support of $cr$.
At a prime dividing $r$, impose depth $8v_p(r)$ on the coefficients of
$f(H)$ and depth $24v_p(r)$ on $z_0,z_1$. Then $W$ is integral there.
At primes dividing $c$ but not $r$, sufficiently positive depths give
integrality as well. The actual pole input $crW^2$ is positive, has
square reduced denominator and the prescribed local squareclasses.
These are the retained pole-completeness hypotheses; $W$ need not be
a unit. Fresh pole witnesses are supplied for this new input.

The scalar coefficient guard may be imposed on $f(R)_1$ instead of
$f(R)_0$. In the reciprocal chart $f(R)=g^2$ one has
$f(R)_0=g_0^2+2g_1^2$ and $f(R)_1=2g_0g_1$. A nonzero second
coefficient therefore implies a positive first coefficient and recovers
the old guard. Completeness uses the established exact formula
$v_3(f([s]P)_1)=2+2v_3(s)$ for nonzero even $s$, so this second
coefficient is nonzero on every prepared scalar orbit point.
This is an arithmetic replacement of the guard, not an equivalence
for arbitrary pairs $g_0,g_1$.

\subsection{The two final height contradictions}
The local order-forty tests force the contact radicals to divide
$F=t^2-a(a-2)tb+(a-2)^2b^2-1$ to order40.
If $F\ne0$, then $|F|<5B^4$ and~\eqref{eq:v13-height} contradicts
the retained lower bound $B\ge|a|>(40\mathcal H)^{10}$.
Thus $F=0$. The nonzero odd root label $b$ and the retained Lucas lemma
give $\log B>|a|-3$.
If $n\ne\beta$, its nonzero reduced numerator $N$ satisfies
$|N|\le2e^{\operatorname h(\beta)}B$ and is divisible by the three
contact radicals to order10. Since $B^2-6>(5/8)B^2$ for $B>4$,
\[
 B^{1/10}<48\mathcal H e^{(19/10)\operatorname h(\beta)},\qquad
 \log B<10\log(48\mathcal H)+19\operatorname h(\beta)
          <29|a|/40<|a|-3.
\]
This contradiction proves $n=\beta$. The unchanged root chart makes $n$ odd, so the parity-sensitive input argument above gives $X\in\mathbb Z$.
The first contradiction uses order40; lowering only the scalar depth
does not weaken that test.

\subsection{Polynomial circuit reductions}
For retained field values $Y_i$, each order block evaluates
$\prod_i(1+qY_i)$. In order1 group the nine factors as
\[
 [V_0,U_1,U_2],\quad[V_1,U_3,U_4],\quad[V_2,V_3,V_4],
\]
where labels denote the squared test values with their existing
conjugations. Order10 uses its pair; order40 retains its pair and
triple. For a triple $a,b,c$ the polynomial is
\[
 1+q(a+b+c)+q^2(ab+ac+bc)+q^3abc.
\]
Its coefficients can be shared among the three detectors. If an old
private coordinate is $d=bc$ and $e_2=a(b+c)+d$, eliminate it by
\[
 d=e_2-a(b+c),\qquad e_3=ae_2-a^2(b+c).
\]
For the chosen first members the square $a^2$ already has a retained
quadratic expression from the logarithm block. Substitution therefore
keeps the new residual quadratic. The reverse map reinserts the displayed
d. Every other use is substituted as well, including affine-normalized
W coordinates. Such a coordinate has the explicit inverse
$W_{\rm old}=(W_{\rm new}-r)/c$ with a fixed nonzero rational c.
These are polynomial reconstructions apart from division by constants.
They do not alter the field test values or pole conditions.

For the order weights let E and q be the retained nonzero scale and
order-one weight. Use
\[
 Ea=q,\quad Eb=a^2,\quad qc=b^2,\quad cd=q^2q^3,
 \quad cq_{10}=d^2,\quad cq_{40}=q_{10}^4.
\]
The last fourth power is evaluated through its retained square.
These equations give $c=q^3/E^6$, $d=q^2E^6$,
$q_{10}=qE^{18}$ and $q_{40}=qE^{78}$. All divisions are licensed by
the unchanged norm/trace guards. The expression $q^5=q^2q^3$
removes the former fifth-power coordinate without increasing degree.

An exact elimination is certified by a residual identity
\[
 F_{\rm new}=F_{\rm old}|_{x=q}
              +\sum_j A_jG_j,
\]
where the retained $G_j$ vanish and the old defining equation has a
nonzero constant pivot in x. Forgetting x proves the forward map;
reinserting x=q proves the reverse map. Checking all consumers is
part of this argument, not an inference from a lower variable count.

\subsection{Factored poles and the preparation conic}
For a normalized detector, the old second-conic rows include
\[
 17E-Eb-17F=0,\quad t_2=t^2,\quad q_2=q^2,\quad
 17bq_2+867F^2-34Ft_2-t_2^2=0.
\]
Replace $F=t v$, $q=t Q$, $q_2=t^2Q_2$ and use
\[
 Q_2=Q^2,\qquad17bQ_2+867v^2-34tv-t^2=0,
 \qquad17E-Eb-17tv=0.
\]
The old conic is $t^2$ times the new conic, so every new solution
reconstructs an old one, including t=0. For completeness the retained
construction supplies positive semantic b. If t!=0 divide by t;
if t=0 the old equation forces F=q=0, so the zero lift works.
The same argument applies to the scalar pole after its exact fixed
normalization. Its stored input has the opposite sign from the semantic
positive input; positivity must not be inferred from the stored symbol.

For scalar preparation, only the square N of the prepared value is
consumed outside its private block. The old equations are
\[
 t_2=t^2,\quad\mathfrak b(T-t_2)=1,\quad TW=2,
 \quad N=(W-\mathfrak b)^2.
\]
Replace them by
\[
 U^2=N,\qquad \mathfrak bJ=1,\qquad JN+Y^2=\mathfrak b.
\]
Forward take $U=W-\mathfrak b$, $J=T-t_2$, $Y=tW$. Conversely, if
$d=\mathfrak b+U\ne0$, take $W=d$, $t=Y/d$, $t_2=t^2$, $T=t_2+J$.
The conic identity gives TW=2. If d=0, then Y=0 and
$N=\mathfrak b^2$; choose $W=2\mathfrak b$, $t=t_2=0$, $T=J$.
This changes only the sign of a private unsquared value and preserves N.
The inverse guard is retained, and no rational point on an unrelated
conic is being assumed.

\subsection{Specializing the scalar preparation to zero}
Write $q=\operatorname{Norm}(f(H))$, and write the actual scalar ratio as
$z(R)/z(H)=z_0+z_1\rho$. The retained expressions are
\[
 B=(z_0^2+rz_1^2)/r^{32},\qquad
 W=(B^2+rq^2)/r^{32},\qquad \mathfrak b=crW^2.
\]
Here $r,c$ are positive and $q\ne0$, so $W,\mathfrak b>0$ at every
rational solution, independently of a private preparation inverse.
The earlier preparation-zero lemma applies when $\mathfrak b<17$, its
reduced denominator is a square, it is locally square at $2,3,17$, and
its valuations there are at least $3,1,2$, respectively. We now arrange
these conditions at the same points used in the final construction.

Let $G$ be the closure of the simultaneous real orbit of $P$ and its
conjugate, let $G^0$ be its identity component, and put $m_0=|G/G^0|$.
The real-progression lemma says that a progression with both residue
and modulus divisible by $m_0$ has dense positive tail in $G^0$.
The functions $f,z$ extend continuously to the identity with value zero.
Choose a multiple $M$ of $64m_0$ whose multiples give all desired formal
depths at primes in $S=\operatorname{supp}(cr)\cup\{2,3,17\}$.
In particular require coefficient depth $8\eta$ for $f(H)$ when
$\eta=v_p(r)>0$, and positive coefficient depth at17. Density permits
$H_0=[h_0]P$, $h_0>0$, $M\mid h_0$, as close to the two real identities
as needed to make $q(H_0)^2/r^{31}$ arbitrarily small.
Its real $z$ values are nonzero.

With $H_0$ fixed, choose a multiple $N_0$ of $6m_0$ imposing the anchor
and all required scalar finite depths. At $p\mid r$ require coefficient
depth $24\eta$ for $z(R)/z(H_0)$; impose positive depths also at17 and
at primes dividing $c$ but not $r$. The formal subgroups have finite
index, so arbitrary further depths can be imposed this way. Density of
the positive progression $s\equiv0\pmod{N_0}$ permits $s$ arbitrarily
large and $R=[s]P$ arbitrarily close to both real identities. The two
real ratios, hence their rational coefficients $z_0,z_1$, become small.
Choose them so that $\mathfrak b(H_0,R)<17/4$. This also permits the
size part of the anchor; set $a=s+1$ as before.

At $p\mid r$ these depths give $v_p(q)\ge16\eta$,
$v_p(B)\ge16\eta$ and $v_p(W)\ge0$. At $p\mid c$, $p\nmid r$,
the imposed positive depths again give $v_p(W)\ge0$. Outside this
support $v_p(\mathfrak b)=2v_p(W)$, so the reduced denominator of
$\mathfrak b$ is a square. Its squareclasses at $2,3,17$ are those of
$cr$. Since $v_2(cr)=4$, $v_3(cr)=2$, the first two required bounds
hold. At17 both $r,c$ are units; the separately imposed depths give
$v_{17}(q),v_{17}(B)\ge2$, hence $v_{17}(W)\ge4$ and
$v_{17}(\mathfrak b)\ge8$. Inclusion of17 does not depend on it
belonging to the support of $cr$.

Keep $R$ fixed. A nonempty real neighborhood $U\subset G^0$ of $H_0$
has $0<\mathfrak b(H,R)<17$. A sufficiently divisible continuity
modulus $L$, divisible by $64m_0$, preserves all the finite conditions
under $h\equiv h_0\pmod L$. Choose the remaining labels, root points
and height scales before the contact primes, and enlarge $L$ for their
finite continuity conditions. The finitely many comparison pole sets
have measure zero in $G^0$ by the real-progression lemma. Choose a
nonempty open $U'$, with compact closure in $U$, avoiding them. All
comparison functions are now uniformly bounded on that closure, by
bounds fixed before any contact prime is chosen.

The original simultaneous-prime and separation argument is valid for
this enlarged fixed $L$. Its compatible congruences are
\[
 h\equiv0\pmod{\Pi^{19}M_-},\qquad
 h\equiv1\pmod{M_+},\qquad h\equiv h_0\pmod L.
\]
Their solution progression has residue and modulus divisible by $m_0$.
Its positive tail is therefore dense in $G^0$ and meets $U'$ at
arbitrarily large indices. Choose the final $H$ there. It has every
required primary and opposite-place depth, all scalar finite depths,
the real comparison bounds and $0<\mathfrak b<17$ simultaneously.
The preparation-zero lemma supplies fresh pole witnesses with prepared
value exactly $\mathfrak b$. A zero factor in the factored pole is
handled by the zero lift already proved above.

For clarity, the physical elimination is also exact. With its fixed
positive normalization $K$, the semantic input is
$\mathfrak b=-Kx$ and the consumed square is $K^2n$. The old rows are
\[
 U^2=K^2n,\qquad1+KxJ=0,\qquad Kx+K^2Jn+V^2=0.
\]
Replace all three by $n=x^2$ and forget $U,J,V$. At any new solution
$\mathfrak b>0$ as proved independently above; restore
$U=\mathfrak b$, $J=1/\mathfrak b$, $V=0$. Each old row is then zero,
and every outside value is unchanged. The opposite direction for the
integer input projection is the simultaneous construction, not a claim
that an arbitrary old preparation witness can be specialized.

\subsection{Normalizing the remaining scalar pole}
After the preparation-zero replacement, let $x$ again denote the
physical scalar input coordinate, so $\mathfrak b=-Kx>0$ and $x\ne0$.
This nonvanishing was proved from the unchanged point and scalar root
rows; it does not depend on the preparation square or on the private
pole equations. Assign weight one to the five private slots
$(e,h,p,q,v)$ and weight two to their square slots $(p_2,q_2)$.
Replace each old slot of weight $j$ by $x^j$ times its new slot,
and discard the preparation-square slot $n$, setting it equal to $x^2$.
Every one of the five remaining pole residuals is then $x^2$ times
its new residual. For example the first conic and the h equation have
respectively their $ex$ and $x^2$ terms replaced by $e$ and1, and
\[
 17ex+Ken-17hv=0
       \quad\longmapsto\quad 17e+Kxe-17hv=0.
\]
The square rows remain $p_2=p^2$, $q_2=q^2$, and the factored
second-conic coefficient list is unchanged. In the physical indexing
of the reviewed component, the weight-one slots are
$51,52,53,54,532$, the weight-two slots are $527,534$, and $531$ is n.
These are implementation indices, not mathematical constants.

From an old solution divide each private slot by its indicated power
of x. From a new solution multiply by that power and restore $n=x^2$.
These maps are inverse and fix every outside coordinate. The sign of x
is immaterial, and no division uses a private pole value. In particular
all zero branches of p,q,v or h are preserved. These seven slots occur only in the five displayed pole rows, and n
has only its former definition and these consumers. Thus there is no
untransformed outside use of a rescaled slot.

\subsection{Rescaling the scalar anchor}
The odd scalar root S has the unchanged chart
$S=(x,hj)$, $x=h^2$, $j^2=x^2-D$, hence $S=[a]P$ for an odd integer a.
Write the retained reciprocal difference equation as
\[
 (c_*x+D)g=a_*j-b_*h,\qquad g=g_0+g_1\rho.
\]
Its denominator is positive in both real embeddings. We replace the
old inverse of the irrational coefficient of f(S-P) by
$g_0\tau=1$. This implies $g\ne0$ before using the anchor.
The subtraction chart identifies $g^2=f(S-P)$: at $S=P$ the quotient
would be zero, which is excluded, and at $S=-P$ the explicit duplication
formula gives the same value $f(-2P)$. Thus $S-P=[s]P$ has a nonzero
even index $s=a-1$ without any prior anchor divisibility assertion.

The earlier scalar coefficient lemma applies to every such s and gives
\[
 v_3(f_1([s]P))=2+2v_3(s).
\]
Its hypotheses contain no anchor equation. Indeed
$z(2P)/3\equiv2+\rho\pmod3$; multiplication by units and then by
powers of3 in the integral formal group preserve the stated leading
term. The curve equation makes the leading residue of
$f/3^{2+2v_3(s)}$ a nonzero rational multiple of $\rho$.
Independence of1 and $\rho$ in the residue field proves the displayed
coefficient valuation, also for negative s. Consequently
$F=f_1([s]P)=2g_0g_1\ne0$, in particular $g_1\ne0$.
This is the noncircular justification for restoring the old inverse.

The old three-row subsystem is
\[
 F=2g_0g_1,\quad F\sigma=1,\quad
 u^2+6v^2-2w^2-3F^2-4Cw=0.
\]
The new subsystem is
\[
 g_0\tau=1,\qquad
 U^2+6V^2-2W^2-12g_1^2-4C\tau W=0.
\]
Forward put $(U,V,W)=(u,v,w)/g_0$, $\tau=1/g_0$ and forget F.
Backward put $(u,v,w)=g_0(U,V,W)$, restore $F=2g_0g_1$ and
$\sigma=1/F$. If $q_{\rm old},q_{\rm new}$ denote the anchor residuals,
the exact identity is
\[
 q_{\rm old}|_{(u,v,w)=g_0(U,V,W),\ F=2g_0g_1}
       -g_0^2q_{\rm new}
       =4Cg_0W(g_0\tau-1).
\]
These inverse maps fix the input and all other coordinates. They permit
zero values of U,V,W. They are equivalences for the full point system,
not for an isolated anchor with arbitrary g: that isolated equation
would admit $g_1=0$, which the coefficient lemma just excluded.

When C remains a fixed integer coefficient, both new residuals are
quadratic and their squared sum contributes $16C^2\tau^2W^2$.
For the ordinary companion compute C by the quadratic chain
$c_0=3^{1024}$, $c_j=c_{j-1}^2$ for $1\le j\le15$, so that
$c_{15}=3^{2^{25}}=C$. Introduce a separate multiplication
slot $T_C$ and equation
\[
                 T_C=C\tau.
\]
Use $-4T_CW$ in the anchor. Treating a computed C as a coefficient would
instead make $C\tau W$ cubic, so this extra slot is necessary for that
companion. Both directions restore $T_C$ uniquely; no nonvanishing or
point-domain argument is changed. This slot and its gate, together with the sixteen constant-chain slots,
are included in the ordinary companion.

\subsection{Completeness and finite certification}
Fix an integer X and compute the positive deterministic scalar correction.
Choose the preliminary shared point and then the deep nonzero scalar
point by the simultaneous real and finite-place construction just proved.
Choose the odd Lucas labels, then their actual root points and positive
root scales. Fix a continuity modulus for all the new rational functions.
Only after these choices select the contact primes and solve the retained
primary/opposite congruences. Take a sufficiently large common index h in the prescribed real open
set, avoiding the finitely many old exceptional point loci. This gives
$U_i=[hn_i]P$ and $V_i=[(h-1)n_i]P$, all nonidentity in this completeness
construction. The new rational logarithms have the required local
precision, and the original real pole estimates and arbitrary finite
preparation depths give positive semantic pole inputs. Obtain fresh
pole witnesses with scalar preparation equal to its input, apply the
factored-conic maps, and reconstruct each
finite private circuit. These are simultaneous choices; no claim is
made that an arbitrary old pole witness survives a changed logarithm
or changed scalar root.

\subsection{The specified polynomials and coefficient tables}
Clear denominators and divide each nonzero quadratic residual by its
integer content. Define $P_{394}$ to be the sum of the squares of the
364 resulting integer residuals. Over the rational numbers this sum
vanishes exactly when every residual vanishes. The soundness and
simultaneous completeness arguments above therefore prove its stated
integer projection. The sixteen constant-chain equations and the
$T_C=C\tau$ equation give the ordinary companion: its 381 normalized
residual squares define $P_{411}$, with the same integer projection.

Expansion of the specified residual lists gives respectively 6,338 and
6,388 distinct nonzero monomials. Both polynomials have total degree
four and content one. For the symbolic polynomial the largest absolute
coefficient is $16C^2=16\cdot3^{2^{26}}$; for the ordinary polynomial it
has 3,247 bits and is strictly less than $2^{3247}$. These counts and
bounds are computed from the complete coefficient lists, not by adding
component savings. Every listed coordinate occurs.

The coefficient artifacts use consecutive coordinate indices, with
index one for X. The symbolic JSON coefficient table has 973,221 bytes,
and the ordinary JSON coefficient table has 982,235 bytes. Their plain
text tables have 212,067 and 214,957 bytes, respectively. The tables,
residual lists, coordinate maps and exact maximum coefficients specify
the two polynomials independently of their factored construction.
